\documentclass[11pt,a4paper]{article}
\usepackage[utf8]{inputenc}
\usepackage[T1]{fontenc}
\usepackage{lmodern,amsmath,amssymb,amsthm,mathtools,booktabs,hyperref,enumitem,microtype}
\usepackage[margin=1in]{geometry}
\usepackage{natbib}
\newtheorem{definition}{Definition}[section]
\newtheorem{proposition}[definition]{Proposition}
\theoremstyle{remark}
\newtheorem{remark}[definition]{Remark}

\title{Causal Attitude Networks on the Substrate\\\large What Moved, What Did Not, and What the Mean Field Costs}
\author{Matthew Parslow\\Independent Researcher}
\date{\textbf{Draft \textemdash\ draft-2026-08-30.4}\\[2pt]1 September 2026 \textbullet\ document revision 1}

\usepackage{fancyhdr}
\pagestyle{fancy}
\fancyhf{}
\fancyfoot[C]{\thepage}
\fancyfoot[R]{\scriptsize\textsc{draft-2026-08-30.4\,r1}}
\renewcommand{\headrulewidth}{0pt}
\fancypagestyle{plain}{\fancyhf{}\fancyfoot[C]{\thepage}%
  \fancyfoot[R]{\scriptsize\textsc{draft-2026-08-30.4\,r1}}\renewcommand{\headrulewidth}{0pt}}
\begin{document}
\maketitle
\begin{center}
\fbox{\parbox{0.92\textwidth}{\small
\textbf{DRAFT --- draft-2026-08-30.4 r1} (1 September 2026).
Circulated for comment, not as a finished result. This paper reports a migration that is
\emph{partial} and says where it stops. Sections marked as falsifiability commitments are the
intended points of attack.}}
\end{center}
\vspace{1em}

\begin{abstract}
The Causal Attitude Network \citep{dalege2016} models an attitude as an Ising-family network of
interacting evaluative components. It was reimplemented on this substrate as a mean-field
self-consistency solver and run. Three things are reported and only one of them is a success. The
\emph{alignment} is incomplete in a specific and locatable way: the dynamics have no term surface in
the substrate's own vocabulary, so the model runs beside the graph rather than in it. The
\emph{performance} question turns out not to be the complexity one it appears to be. And of two
registered predictions, one is confirmed with a large effect and a control, and one is
\emph{reversed} --- reported as measured rather than reinterpreted.
\end{abstract}

\section{What was moved, and what a move would mean}\label{sec:moved}
An attitude in the CAN model is a network of evaluative components with pairwise couplings and
external fields, and its behaviour is the behaviour of an Ising system: relaxation toward low-energy
configurations, frustration where couplings conflict, and a temperature-like parameter governing how
sharply it settles. The reimplementation supplies that as a solver on the graph, exposed to the host
and driven by seed-resident network data.

\emph{Moving a formalism onto this substrate means more than making it run there.} The substrate's
claim is that decisions are taken by graph rules over graph data; a model that computes its result in
a compiled routine and hands back an answer has been \emph{hosted}, not moved. The distinction is the
whole of the alignment question below, and it is not rhetorical: the two cases differ in what can
subsequently be inspected, mutated, minted and reasoned over by the rest of the account.

\section{Alignment: the dynamics have no term surface}\label{sec:alignment}
Fifteen rules written against the network were inert. The obvious diagnosis --- missing or malformed
match clauses --- is not the reason.

\emph{There is no DSL term surface for the dynamics at all.} The relaxation exists only across a
host-language boundary, so a rule cannot express it, cannot be matched against it, and cannot mint
against it. The rules were not broken; there was nothing for them to be written in terms of. That is
a structural finding about the migration rather than a defect in fifteen rules, and it is why this
paper describes the move as partial.

One rule was then wired with real match, where and effects clauses using only the existing term
vocabulary, and works on the production path. \emph{So the boundary is not a wall}, and the correct
statement is the narrow one: what has been moved is the network's data and its results, not its
dynamics, and moving the dynamics is available work rather than blocked work.

\section{Performance: the complexity question is the wrong one}\label{sec:performance}
The natural objection to putting an Ising model anywhere is that its ground-state problem is hard.
\citet{barahona1982} makes the sharp version available: the planar two-dimensional problem
\emph{without} an external field is solvable in polynomial time, while the same planar problem
\emph{with} a field, and the three-dimensional problem either way, is NP-hard. Every network here
carries a nonzero field, so the hard case is the relevant one.

\emph{And it does not bind, because no exact ground state is being sought.} The solver computes a
mean-field self-consistency fixed point --- polynomial per iteration, unconditionally, at every
network size. The actionable question is therefore not whether this scales exactly, since no
algorithm does and none is attempted, but how good the approximation is and when it fails to
converge at all.

\begin{remark}[The weakest member of a hierarchy that was borrowed whole]\label{rem:hierarchy}
Mean field is the first and weakest step of a known ladder --- mean field, then Bethe and loopy
belief propagation, then Kikuchi and generalized belief propagation \citep{yedidia2005} --- and it is
weakest \emph{precisely} in the regime of loops, strong local correlation and frustration. That is
the regime this implementation's own frustration and basin audits exist to detect, so the
approximation is thinnest exactly where the measurement is aimed.

The honest form of that observation is not that a mistake was made. \emph{An entire coherent piece of
statistical-mechanics vocabulary was borrowed, and the better half of the same literature is sitting
one step away unused.} A caveat comes with it and is not decoration: Bethe and loopy propagation are
themselves known to lose stable solutions in the same frustrated, field-bearing, low-temperature
regime, so the ladder's next rung is an improvement and not a cure.
\end{remark}

\section{Results}\label{sec:results}

\begin{remark}[The registered scaling prediction is reversed]\label{rem:null}
The primary registered prediction --- that coarse quantization would yield fewer distinct basins than
fine, with the gap widening in network size --- \emph{is not confirmed, and the measured direction is
closer to reversed}. Coarse yields \emph{more} distinct basins ($6$ against $4$ at $N=4$; $15$
against $4$ at $N=8$), a higher entropy-weighted effective basin count, and higher mean frustration
($2.10$ against $1.55$ at $N=4$; $15.32$ against $11.06$ at $N=8$). At $N=16$ the two conditions
converge to near-parity on every measure.

The manipulation check holds throughout --- coarse quantization does produce larger per-edge
representational error at every $N$ --- so the intervention worked and the prediction about its
consequence was wrong. That is the useful shape of a null and it is reported here rather than
reinterpreted into a different prediction that would have survived.
\end{remark}

\begin{remark}[Resolution separates the basins, with a control that says so]\label{rem:bigville}
A second registered run, on a population of simulated residents each carrying its own network at its
own representational resolution, confirmed both of its predictions.

The manipulation check is large: the coarse resident's mean rounding error is about $9.2$ times the
fine resident's ($0.553$ against $0.060$). The primary prediction --- deliberately not the naive
basin-count intuition, that having already been falsified above on the same engine --- is that under
\emph{identical} scenario evidence the coarse network relaxes to a positive-energy, highly frustrated
basin (frustration $9.86$, energy $+5.19$) while the fine one relaxes to a negative-energy,
low-frustration basin (frustration $1.62$, energy $-1.77$). It does.

A population-wide check across one hundred residents, reported as exploratory rather than
pre-registered to the same standard, shows the same direction on average (mean frustration $3.29$ in
the low-resolution group against $1.38$ in the high). \emph{The control is what makes this worth
reporting}: holding resolution fixed across the population collapses the trait-to-frustration
correlation from $-0.33$ to $-0.08$, so resolution rather than some correlated trait is carrying it.
\end{remark}

\begin{remark}[A discriminative null was found, not built]\label{rem:found}
The relaxation exhibits a genuine binary between cohering and failing to cohere: clean convergence in
one regime, a persistent two-cycle oscillation in another, confirmed against a same-parity control.

That shape --- an operation that either resolves or returns a determinate nothing --- is one this
programme has elsewhere and did not put here. \emph{It was already present, unexploited, in the
borrowed model's own mathematics.} What is \emph{not} established is the tempting next step: whether
low resolution specifically drives the failure, as against the particular realised coupling draw, is
a weak and noisy same-direction lean across two independent samples of $24$ and $40$ trials. It is
reported as a lean and not as an effect.
\end{remark}

\section{Falsifiability commitments}\label{sec:falsifiability}
\begin{enumerate}[leftmargin=*]
\item \textsc{[shape]}~\textbf{The migration is partial in the way stated.} Fails if the dynamics can
be expressed in the existing term vocabulary without new term surface --- that is, if the fifteen
rules were inert for an ordinary reason after all.
\item \textsc{[detail]}~\textbf{Resolution drives basin quality.} Fails if the frustration and energy
separation survives holding resolution fixed, or if the population correlation does not collapse
under that control.
\item \textsc{[detail]}~\textbf{The scaling prediction is reversed and not merely unconfirmed.} Fails
if the basin counts recover the registered direction under a larger sample at $N=4$ and $N=8$.
\item \textsc{[shape]}~\textbf{Mean field is the binding limitation, not complexity.} Fails if an
exact solver at these network sizes changes the reported quantities materially, which would mean the
approximation rather than the model was being measured.
\item \textsc{[shape]}~\textbf{The oscillation is the model's and not the implementation's.} Fails if
the two-cycle disappears under an equivalent solver written independently.
\end{enumerate}

\begin{thebibliography}{9}

\bibitem[Barahona(1982)]{barahona1982}
Francisco Barahona.
\newblock On the computational complexity of Ising spin glass models.
\newblock \emph{Journal of Physics A: Mathematical and General}, 15(10):3241--3253, 1982.

\bibitem[Dalege et al.(2016)]{dalege2016}
Jonas Dalege, Denny Borsboom, Frenk van Harreveld, Helma van den Berg, Mark Conner, and Han L.~J.
van der Maas.
\newblock Toward a formalized account of attitudes: The Causal Attitude Network (CAN) model.
\newblock \emph{Psychological Review}, 123(1):2--22, 2016.

\bibitem[Yedidia et al.(2005)]{yedidia2005}
Jonathan~S. Yedidia, William~T. Freeman, and Yair Weiss.
\newblock Constructing free-energy approximations and generalized belief propagation algorithms.
\newblock \emph{IEEE Transactions on Information Theory}, 51(7):2282--2312, 2005.

\end{thebibliography}
\end{document}
